\uuid{Lq8D}
\titre{Compter les paramètres d'un RNN, LSTM et GRU}
\niveau{L3}
\module{Apprentissage automatique}
\chapitre{Réseaux de neurones récurrents}
\sousChapitre{Architecture et complexité}
\theme{Nombre de paramètres, RNN, LSTM, GRU, capacité du modèle}
\auteur{Maxime Nguyen}
\datecreate{2026-05-13}
\organisation{AMSCC}
\difficulte{1}

\contenu{
\texte{
  Le nombre de paramètres d'un réseau est un indicateur clé de sa capacité et de son coût d'entraînement. On note :
  \begin{itemize}
    \item $d$ : dimension de l'entrée $x_t$,
    \item $h$ : dimension de l'état caché,
    \item $p$ : dimension de la sortie $y_t$.
  \end{itemize}
  Chaque architecture possède une \textbf{couche de sortie linéaire} $y_t = W_y h_t + b_y$.

  \medskip
  \textbf{Rappel des architectures.}

  \textit{RNN simple :}
  \[
    h_t = \tanh(W_x x_t + W_h h_{t-1} + b_h)
  \]

  \textit{LSTM} — quatre portes, chacune de la forme $W x_t + U h_{t-1} + b$ :
  \[
    f_t, \quad i_t, \quad \tilde{c}_t, \quad o_t
  \]

  \textit{GRU} — trois portes :
  \begin{align*}
    r_t &= \sigma(W_r x_t + U_r h_{t-1} + b_r) & &\text{(porte de réinitialisation)}\\
    z_t &= \sigma(W_z x_t + U_z h_{t-1} + b_z) & &\text{(porte de mise à jour)}\\
    \tilde{h}_t &= \tanh(W_h x_t + U_h (r_t \odot h_{t-1}) + b_h) & &\text{(candidat)}\\
    h_t &= (1 - z_t) \odot h_{t-1} + z_t \odot \tilde{h}_t
  \end{align*}
}
\begin{enumerate}
\item
  \question{Justifier que le nombre de paramètres du RNN simple (hors couche de sortie) est $N_{\text{RNN}} = h(d + h + 1)$. Détailler la contribution de chaque matrice et biais.}
  \reponse{
    Le RNN simple possède trois jeux de paramètres :
    \begin{itemize}
      \item $W_x \in \mathbb{R}^{h \times d}$ : $hd$ paramètres,
      \item $W_h \in \mathbb{R}^{h \times h}$ : $h^2$ paramètres,
      \item $b_h \in \mathbb{R}^h$ : $h$ paramètres.
    \end{itemize}
    Total : $hd + h^2 + h = h(d + h + 1)$.
  }

\item
  \question{Montrer que le nombre de paramètres de la LSTM (hors couche de sortie) est $N_{\text{LSTM}} = 4h(d + h + 1)$. Pourquoi ce facteur $4$ ?}
  \reponse{
    Chacune des quatre portes ($f$, $i$, $\tilde{c}$, $o$) possède exactement la même structure qu'un RNN simple : une matrice $W \in \mathbb{R}^{h \times d}$, une matrice $U \in \mathbb{R}^{h \times h}$ et un biais $b \in \mathbb{R}^h$, soit $h(d + h + 1)$ paramètres par porte. Le facteur $4$ provient du fait qu'il y a quatre portes indépendantes :
    \[
      N_{\text{LSTM}} = 4 \times h(d + h + 1).
    \]
  }

\item
  \question{Montrer que le nombre de paramètres du GRU (hors couche de sortie) est $N_{\text{GRU}} = 3h(d + h + 1)$.}
  \reponse{
    Le GRU possède trois portes ($r$, $z$, $\tilde{h}$), chacune avec la même structure : $W \in \mathbb{R}^{h \times d}$, $U \in \mathbb{R}^{h \times h}$ et $b \in \mathbb{R}^h$. (Remarque : dans $\tilde{h}_t$, $U_h$ multiplie $r_t \odot h_{t-1}$, mais $U_h$ reste de taille $h \times h$.) D'où :
    \[
      N_{\text{GRU}} = 3 \times h(d + h + 1).
    \]
  }

\item
  \question{Ajouter la contribution de $W_y \in \mathbb{R}^{p \times h}$ et $b_y \in \mathbb{R}^p$. Écrire le nombre total de paramètres $N_{\text{total}}$ pour chacune des trois architectures.}
  \reponse{
    La couche de sortie ajoute $ph + p = p(h+1)$ paramètres dans les trois cas.
    \begin{align*}
      N_{\text{total}}^{\text{RNN}}  &= h(d + h + 1) + p(h + 1) \\
      N_{\text{total}}^{\text{LSTM}} &= 4h(d + h + 1) + p(h + 1) \\
      N_{\text{total}}^{\text{GRU}}  &= 3h(d + h + 1) + p(h + 1)
    \end{align*}
  }

\item
  \question{Application numérique : $d = 10$, $h = 64$, $p = 5$. Compléter le tableau suivant.

  \begin{center}
  \begin{tabular}{lcc}
    \hline
    Architecture & Paramètres (hors sortie) & Paramètres (avec sortie) \\
    \hline
    RNN simple & & \\
    GRU        & & \\
    LSTM       & & \\
    \hline
  \end{tabular}
  \end{center}
  }
  \reponse{
    $h(d + h + 1) = 64 \times (10 + 64 + 1) = 64 \times 75 = 4\,800$.

    $p(h + 1) = 5 \times 65 = 325$.

    \begin{center}
    \begin{tabular}{lcc}
      \hline
      Architecture & Hors sortie & Avec sortie \\
      \hline
      RNN simple & $4\,800$  & $5\,125$  \\
      GRU        & $14\,400$ & $14\,725$ \\
      LSTM       & $19\,200$ & $19\,525$ \\
      \hline
    \end{tabular}
    \end{center}
  }

\item
  \question{La LSTM a davantage de paramètres que le GRU, mais est-elle toujours plus performante ? Dans quel cas préférerait-on le GRU ?}
  \reponse{
    Non, une plus grande capacité ne garantit pas de meilleures performances. Le GRU est souvent préféré dans les situations suivantes :
    \begin{itemize}
      \item \textbf{Données peu abondantes} : moins de paramètres réduit le risque de surapprentissage.
      \item \textbf{Contrainte de temps ou de mémoire} : le GRU est environ $25\,\%$ plus rapide à entraîner que la LSTM pour une même taille de couche cachée.
      \item \textbf{Séquences de longueur modérée} : la LSTM apporte un avantage surtout pour des dépendances très longues ; sur des séquences courtes, les deux architectures sont comparables.
    \end{itemize}
    En pratique, le choix se fait souvent par validation croisée sur la tâche cible.
  }
\end{enumerate}
}
